%% Revere: A Structured Workflow for Auditable Political Reasoning
%% Candidacy Project Technical Writeup
%%
%% Compile: pdflatex revere_system_design.tex  (run twice for cross-references)

\documentclass[11pt,letterpaper]{article}

\usepackage[margin=1.1in]{geometry}
\usepackage[T1]{fontenc}
\usepackage{xcolor}
\usepackage{enumitem}
\usepackage[colorlinks=true, linkcolor=blue!60!black, urlcolor=blue!60!black,
            citecolor=blue!60!black]{hyperref}
\usepackage{array}
\usepackage{booktabs}
\usepackage{parskip}
\usepackage{listings}
\usepackage{microtype}

\definecolor{muted}{RGB}{100,100,100}
\definecolor{todo}{RGB}{160,40,40}
\definecolor{codebg}{RGB}{245,245,245}

\lstset{
  basicstyle=\ttfamily\small,
  backgroundcolor=\color{codebg},
  frame=single, framerule=0.4pt, rulecolor=\color{gray!40},
  breaklines=true, captionpos=b, aboveskip=6pt, belowskip=6pt, numbers=none,
}

\newcommand{\todo}[1]{\textcolor{todo}{\textbf{[TODO: #1]}}}
\newcommand{\kw}[1]{\textbf{#1}}
\newcommand{\code}[1]{\texttt{\small #1}}
\newcommand{\file}[1]{\texttt{\small #1}}
\newcommand{\provisional}[1]{\textcolor{muted}{\textit{#1}}}

\setlist[itemize]{leftmargin=1.5em, itemsep=2pt, topsep=4pt}
\setlist[enumerate]{leftmargin=1.5em, itemsep=2pt, topsep=4pt}
\setlist[description]{leftmargin=1.5em, itemsep=3pt, topsep=4pt}

\title{%
  \textbf{Revere: A Structured Workflow for\\[3pt]Auditable Political Reasoning}\\[8pt]
  \large Candidacy Project Technical Writeup
}
\author{Daniel Enriquez\\
  New York University\\
  \href{mailto:dee4002@nyu.edu}{dee4002@nyu.edu}}
\date{\provisional{Draft — June 2026. Candidacy Project Technical Writeup.}}

\begin{document}
\maketitle

\begin{abstract}
Revere is a political-reasoning agent whose central design choice is an
explicit, sequential, seven-stage workflow. Each stage makes one structured
LLM call and produces a Pydantic-validated output that becomes the typed input
for the next stage. The orchestrator owns all control flow; the model owns all
semantic reasoning within each stage's contract.

This separation produces a readable record of every reasoning decision ---
scope classification, search planning, per-source quality assessment,
multi-perspective synthesis, claim verification, and a seven-point neutrality
self-check --- without depending on raw chain-of-thought. A deterministic
\code{EvidenceSufficiency} layer propagates structured coverage flags from
Stage~4 into later stages, enforcing evidence-proportional hedging without
additional model calls. Scope decisions are made by reasoning against a written
charter, not by matching words in the query. This document describes the design
rationale, stage contracts, evaluation methodology, and known limitations.
\end{abstract}


%% ─────────────────────────────────────────────────────────────────────────────
\section{Introduction}
\label{sec:intro}

Political AI systems face a credibility problem. A model that answers
political questions fluently is not therefore answering them fairly or
accurately. The failure mode is not hallucination in the narrow sense ---
plausible-sounding words --- but authority fabrication: presenting contested
claims with the same confidence as established facts, or distributing
analytical effort asymmetrically across perspectives without acknowledging
that asymmetry.

Revere addresses this by making reasoning inspectable. Every query triggers a
fixed sequence of typed reasoning steps, and the user can follow each one: which
sources were assessed and how, what coverage gaps the evidence had, which
specific claims were verified and to what confidence, and whether the final
answer passed a structured neutrality check. The goal is not to eliminate model
uncertainty --- that is impossible --- but to keep the \kw{source of truth
visible}: the citations, source assessments, and claim-verification records
remain readable independently of the prose answer.

For a political assistant, these properties are not cosmetic. A system that
classifies boundary cases by keyword matching cannot explain why it refused a
question; a system that hedges claims uniformly cannot signal which ones are
actually in doubt. Revere delegates all such judgments to typed LLM stage
outputs, keeping control flow in deterministic Python. The rest of this document
describes how that separation is achieved and what it costs.


%% ─────────────────────────────────────────────────────────────────────────────
\section{Design Goals and Non-Goals}
\label{sec:goals}

The system has five design goals. \kw{Auditable reasoning}: every inference
step should be visible as a typed structured output, independently inspectable
from the final answer. \kw{Structured workflow, not free-form agent}: the
pipeline is a deterministic sequence with no planning loop, no tool-calling
cycle, and no self-directing decomposition; the orchestrator decides what runs,
the model decides what it means. \kw{Scope reasoning, not keyword filtering}:
boundary decisions are produced by LLM reasoning against a written charter;
matching words in the query is a design pattern the system deliberately avoids.
\kw{Evidence-proportional neutrality}: perspectives are represented in
proportion to their retrieved evidence base, applying \kw{equal interpretive
respect, not equal factual authority} --- weakly sourced views are attributed
rather than asserted. \kw{Source-of-truth visibility}: citations, per-source
assessments, coverage flags, and verification records are first-class outputs,
not internal states.

Three things are out of scope. Search is used for freshness and grounding, not
as a definitive authority; the system does not do real-time fact-checking.
Neutrality is about interpretive depth and sourcing quality, not word counts.
The audit trace contains structured stage outputs, not the model's private
chain-of-thought; those tokens are never captured.


%% ─────────────────────────────────────────────────────────────────────────────
\section{System Architecture}
\label{sec:arch}

\subsection{Layers}

Revere has four layers. The \kw{schema layer} (\file{revere\_agent/schemas/})
defines Pydantic models for every stage input and output, all using
\code{ConfigDict(extra=\char`"forbid\char`")} so that extra fields from the
model fail validation rather than passing silently. The \kw{stage layer}
(\file{revere\_agent/agent/stages/}) contains one Python function per stage;
each assembles a prompt from prior-stage outputs, calls the LLM via a shared
helper, and returns a validated Pydantic object. No stage has side effects
beyond its return value. The \kw{orchestrator layer}
(\file{revere\_agent/agent/orchestrator.py}) sequences stage calls, owns the
one deterministic branch (the out-of-scope short-circuit after S2), runs
parallel search queries, derives \code{EvidenceSufficiency} from S4 output,
and assembles the reasoning trace. The \kw{provider layer}
(\file{revere\_agent/llm/}) is a Python Protocol abstraction with two concrete
backends: \code{AnthropicProvider} (direct API, used in development) and
\code{BedrockProvider} (AWS Bedrock, the production target). Stage code refers
only to logical aliases (\code{sonnet-main}, \code{haiku-fast}); each provider
resolves those internally, so switching backends is an environment-variable
change.

The FastAPI server wraps the orchestrator in a Server-Sent Events endpoint,
streaming progress updates as each stage completes. The Next.js frontend
proxies the stream without buffering and renders a live activity trace. S1 and
S3 run on \code{haiku-fast} because they perform bounded structured operations
(query normalization, search-query formulation) that require no political
judgment. S2, S4, S5, S6, and S7 run on \code{sonnet-main}; scope decisions
in particular are kept on the stronger model because a mistaken out-of-scope
rejection is a hard failure.

\subsection{Pipeline Diagram}

\begin{center}
\begin{verbatim}
User message
     |
  [S1] Intake ----------------------- IntakeAnalysis
     |                                (canonical query, modality,
     |                                 multi_turn_dependency)
  [S2] Scope  ----------------------- ScopeDecision
     |   |                            (in_scope, reasoning,
     |   |                             confidence, redirect)
     |   +-- in_scope=False --> short-circuit
     |         (S3-S7 do not run)
     |
  [S3] Search Plan ------------------ SearchPlan
     |   |                            (needs_search, queries <=3,
     |   |                             target_source_types)
     |   +-- needs_search=True --> Tavily API (concurrent, <=4 workers,
     |                                         deduped, <=6 hits passed on)
     |
  [S4] Source Quality --------------- EvidenceBase + EvidenceCoverage
     |                                (assessments, confidence,
     |                                 gaps, structured coverage flags)
     |
  [deterministic] EvidenceSufficiency (no LLM call)
     |   derived from S4 coverage + S3 plan
     |   propagated as binding constraint note
     |   into S5, S6, S7 prompts
     |
  [S5] Perspectives ----------------- PerspectiveAnalysis
     |                                (labelled views, core_claims,
     |                                 evidence_coverage rating)
     |
  [S6] Verification ----------------- VerificationReport
     |                                (FactualClaim list with claim_type,
     |                                 support_level, drop_if_uncorroborated)
     |
  [S7] Compose + Check -------------- FinalResponse
                                       (response_text, citations,
                                        SelfCheckReport 7-criteria)
\end{verbatim}
\end{center}


%% ─────────────────────────────────────────────────────────────────────────────
\section{Stage-by-Stage Workflow}
\label{sec:stages}

\begin{center}
\renewcommand{\arraystretch}{1.2}
\begin{tabular}{c p{2.5cm} p{2.8cm} p{5.3cm}}
\toprule
\textbf{Stage} & \textbf{Input} & \textbf{Output} & \textbf{Responsibility} \\
\midrule
S1 & User message + history & \code{IntakeAnalysis} &
  Normalise query; resolve cross-turn referents; classify modality \\
S2 & \code{IntakeAnalysis} & \code{ScopeDecision} &
  Reason against charter; produce \code{in\_scope} and reasoning \\
S3 & \code{IntakeAnalysis} & \code{SearchPlan} &
  Decide whether retrieval adds value; formulate $\leq$3 queries \\
Search & \code{SearchPlan} & \code{list[SearchHit]} &
  Tavily API, concurrent execution, URL dedup (deterministic) \\
S4 & Hits + \code{IntakeAnalysis} & \code{EvidenceBase} &
  Per-source quality assessment; emit \code{EvidenceCoverage} \\
Sufficiency & \code{EvidenceBase} + \code{SearchPlan} & \code{EvidenceSufficiency} &
  Derive coverage flags (deterministic, no LLM) \\
S5 & \code{EvidenceBase} + sufficiency & \code{PerspectiveAnalysis} &
  Steelman each perspective; rate evidence coverage \\
S6 & \code{EvidenceBase} + perspectives + sufficiency & \code{VerificationReport} &
  Decompose claims; tag type, support level, do-not-assert \\
S7 & All prior outputs & \code{FinalResponse} &
  Draft answer; run 7-point \code{SelfCheckReport}; emit citations \\
\bottomrule
\end{tabular}
\end{center}

\bigskip

\subsection{S1 --- Intake and Context Normalisation}

S1 converts the raw user message and conversation history into a
\code{canonical\_query}: a self-contained, pronoun-resolved restatement that
all downstream stages receive as their point of reference. It also classifies
the query's modality and sets a \code{multi\_turn\_dependency} flag. Without
this stage, multi-turn sessions would require each later stage to re-parse
history independently, creating inconsistency. Because S1 performs NLP
normalization with no political judgment, it runs on the smaller model.

\subsection{S2 --- Scope Reasoning}

S2 decides whether the query falls within the agent's political-events domain
by reasoning against a written system charter, not by matching words. The
\code{ScopeDecision} schema's \code{reasoning} field requires the model to
cite charter clauses rather than vocabulary, and the prompt opens with the
explicit instruction: \textit{You MUST NOT use keyword matching.} If the
decision is out-of-scope, the orchestrator short-circuits: S3 through S7 do
not run, and \code{suggested\_redirect} from \code{ScopeDecision} is returned
directly as the boundary response. This makes boundary handling cheap --- two
LLM calls --- while leaving the reasoning that produced it in the trace.

\subsection{S3 --- Search Planning}

S3 decides whether external retrieval adds value for the current query, emitting
a \code{SearchPlan} with a \code{needs\_search} flag, up to three queries, and
a list of target source types. When search is not needed, execution is skipped
and S4 runs in a no-search mode that documents the model-knowledge limitation
in the evidence base. S3 runs on the smaller model because the output is
bounded: a binary decision and a short query list.

\subsection{Search Execution}

Search execution is deterministic I/O with no LLM call. All queries run
concurrently via a thread pool (max four workers); Python's \code{executor.map}
preserves query order, making downstream URL deduplication deterministic
regardless of network timing. At most six unique hits are passed to S4.

\subsection{S4 --- Source Quality Reasoning}

S4 judges each retrieved source on its type, authority, editorial slant,
recency, and relevance, producing per-source \code{SourceAssessment} records
and an aggregate confidence score. When no search ran, it produces a
lightweight record documenting the model-knowledge limitation rather than
fabricating assessments. The key output is a structured
\code{EvidenceCoverage} object --- four typed booleans encoding whether the
evidence base includes primary, court, or government sources, and whether
coverage is asymmetric across perspectives. These booleans encode the semantic
judgments that S4 is positioned to make, so that the next step can consume
them without re-interpreting URL strings or prose text.

\subsection{EvidenceSufficiency (deterministic)}

\code{EvidenceSufficiency} is derived deterministically from
\code{EvidenceCoverage} and the source types S3 requested --- no LLM call,
no heuristics. It sets flags for low-confidence evidence, missing primary
sources, unmet source-type requests, and asymmetric perspective coverage. When
any flag is active, a compact binding-constraint note is injected into the
S5, S6, and S7 prompts, requiring those stages to attribute or hedge claims
from underserved perspectives rather than state them declaratively.

\subsection{S5 --- Multi-Perspective Synthesis}

S5 identifies the major political perspectives on the query and steelmans each
one: framing it as a proponent would, not as an opponent would characterize it.
Each perspective receives an \code{evidence\_coverage} rating (ranging from
\texttt{sourced} to \texttt{none}) that the composition stage uses to calibrate
language. Perspectives with thin source support must use attribution-first
phrasing; perspectives well-grounded in the retrieved evidence may be stated
more directly. This is the principle of equal interpretive respect, not equal
factual authority.

\subsection{S6 --- Claim Verification and Calibration}

Before composing prose, S6 decomposes the planned response into atomic factual
claims and rates each for confidence, verification basis, and support level.
Claims flagged \code{drop\_if\_uncorroborated} must be removed by S7, not
hedged. S6 also produces a \code{things\_i\_should\_not\_assert} list that
S7 carries into composition. Separating this verification pass from composition
is the core idea of Chain-of-Verification (Dhuliawala et al., 2023): calibrated
hedging requires knowing which claims are actually in doubt before writing.

\subsection{S7 --- Composition and Neutrality Self-Check}

S7 writes the user-facing answer, applying S6's do-not-assert rules and the
\code{EvidenceSufficiency} constraints. After composing a draft, it runs a
\code{SelfCheckReport}: seven boolean criteria derived from Anthropic's
political-evenhandedness principles (November 2025), covering whether the
response avoids unsolicited opinion, steelmans each perspective, uses neutral
terminology, distributes depth equally, and weights empirical claims
proportionally to their source support. Revisions made to satisfy any failing
criterion are recorded in the report alongside the final text. S7 compacts its
inputs aggressively before the LLM call, reducing S4--S6 verbose reasoning
fields to the structured conclusions S7 actually needs.


%% ─────────────────────────────────────────────────────────────────────────────
\section{Prompt Contracts as Interfaces}
\label{sec:prompts}

Each stage prompt in \file{revere\_agent/prompts/} acts as an interface
contract: it specifies which schema fields to populate, what semantic task to
perform, and what constraints apply. The system charter is prepended to every
stage's system message, so the same political-neutrality principles are in
context for every LLM call. A 12-character SHA1 prefix of each prompt file is
computed at load time and stored in the trace entry as \code{prompt\_version},
linking any observed behavior change to a specific prompt edit without relying
on git history alone.

Three design choices are worth noting. S2's prompt opens with the literal
instruction \textit{You MUST NOT use keyword matching}, and the schema's
\code{reasoning} field requires charter clause references --- the prompt and
schema reinforce the same constraint. S5's steelmanning requirement is framed
as an Ideological Turing Test: would a proponent of this view endorse the
framing? A self-assessment field in the schema asks the model to answer that
question explicitly. S7's \code{SelfCheckReport} criteria were derived from
Anthropic's published political-evenhandedness measurement work (November
2025), so each boolean maps to a specific property that independent evaluations
can also measure.


%% ─────────────────────────────────────────────────────────────────────────────
\section{Source-of-Truth and Evidence Sufficiency}
\label{sec:sourceoftruth}

A recurring failure mode in AI-assisted research is \kw{authority fabrication}:
the system cites a source that supports its conclusion, but the user has no way
to assess whether the source is real, authoritative, or accurately paraphrased.
Revere's response is to make the source base inspectable. The S4 per-source
assessments, the \code{EvidenceCoverage} flags, and the S6 claim-verification
records are all structured outputs that a reader can examine independently of
the prose answer. Citations are typed objects in the final response schema, not
incidental footnotes.

The design of \code{EvidenceSufficiency} illustrates why this matters
structurally. An earlier version computed coverage sufficiency in deterministic
Python by (a) expanding source types based on URL domain patterns
(\code{supremecourt.gov} $\to$ \texttt{court}; any \code{.gov} TLD $\to$
\texttt{government}) and (b) scanning S4's free-text gap descriptions for
phrases like \textit{one-sided} or \textit{asymmetric}. Both are semantic
inference in deterministic code: they extract meaning from strings and
LLM-generated prose outside of a structured LLM call, making the inference
invisible to evaluation.

The current design eliminates both heuristics. S4 now emits a structured
\code{EvidenceCoverage} object --- four typed booleans and an optional
\code{asymmetry\_note} --- encoding the coverage judgments that S4 is
positioned to make. The deterministic \code{derive\_evidence\_sufficiency}
function reads those booleans directly. No URL parsing, no keyword scanning.
When \code{EvidenceCoverage} is absent (a fallback for legacy sessions),
the function uses source-type enum comparison only and skips asymmetry
detection entirely.

\begin{lstlisting}[caption={\code{EvidenceCoverage} schema
(\file{revere\_agent/schemas/source\_quality.py})}]
class EvidenceCoverage(BaseModel):
    has_primary_sources: bool
    has_court_sources: bool
    has_government_sources: bool
    perspective_coverage_asymmetric: bool
    asymmetry_note: str | None = None
\end{lstlisting}


%% ─────────────────────────────────────────────────────────────────────────────
\section{Structured Audit Trace and User Interface}
\label{sec:ui}

\subsection{Trace Design}

After each stage runs, the orchestrator appends a \code{StageTraceEntry} to
the \code{ReasoningTrace}. Each entry records the stage identifier, a SHA1
prefix of the prompt file used, the model alias, the wall-clock duration, and
the stage's Pydantic output serialised as a plain dictionary. Keeping the
output as a dictionary makes the trace schema-agnostic: it does not need to
enumerate the output types of every stage it might contain. The trace is a
first-class deliverable, not a debug log. A reader inspecting it sees, for
example, the \code{ScopeDecision.reasoning} field --- an LLM-produced string
within a typed schema --- rather than raw model output tokens. This is the
distinction between a structured audit trace and raw chain-of-thought.

\subsection{User Interface}

The Next.js frontend surfaces the same structured stage outputs the backend
produces. A live activity trace updates as each stage completes, showing the
stage name, duration, and a one-line summary. Selecting a stage opens a detail
panel that renders its structured output: the scope reasoning text, per-source
assessments, perspective labels with coverage ratings, factual-claim confidence
cards, and the final \code{SelfCheckReport}. No reasoning logic is implemented
in TypeScript; the panel is a view over the Python pipeline's typed outputs.
A Gradio interface (\file{revere\_agent/ui/gradio\_app.py}) provides an
alternative local entry point for quick testing without the Next.js server.


%% ─────────────────────────────────────────────────────────────────────────────
\section{Evaluation Methodology}
\label{sec:eval}

\subsection{Structural Tests}

The test suite contains 163 tests across 12 files, all using a
\code{StubLLMProvider} that returns pre-configured Pydantic objects without
calling the API. These tests verify pipeline correctness: stage contracts,
the S2 short-circuit, concurrent search ordering, \code{EvidenceSufficiency}
flag derivation, and S7 output compaction. Structural tests are regression
safety, not behavioral quality measurement. They confirm the pipeline behaves
as specified; they do not assess whether any particular response is politically
neutral.

\subsection{Behavioral Evaluation}

Behavioral quality is assessed by running the five demonstration scenarios
(Section~\ref{app:scenarios}) and inspecting the pipeline trace for each. The
relevant signals are: the S2 scope decision and its reasoning text; which
source types S3 requested; whether S4 found those types; the S5 perspective
labels and evidence-coverage ratings; the S6 claim types and any
do-not-assert decisions; and whether all seven S7 \code{SelfCheckReport}
criteria pass. A planned paired-prompt protocol --- running the same query
with neutrality-biasing framing and comparing perspective-coverage outcomes
--- would provide independent evidence beyond the model's self-report, but
has not yet been implemented.


%% ─────────────────────────────────────────────────────────────────────────────
\section{Demonstration Cases and Results}
\label{sec:demos}

\provisional{This section will be filled with observed results after the demo
recording. The five scenarios, in order:}

\begin{enumerate}
  \item \textit{2023 debt ceiling negotiations} --- both parties' positions
        and the final agreement.
  \item \textit{2024 presidential primary issues} --- information freshness
        and balanced perspective coverage.
  \item \textit{SFFA v.\ Harvard and UNC} --- affirmative action ruling;
        legal, political, and social perspectives.
  \item \textit{Weather / homework boundary management} --- graceful
        out-of-scope handling via S2 reasoning.
  \item \textit{Current immigration policy debate} --- perspective balance
        and evidence-proportional language.
\end{enumerate}


%% ─────────────────────────────────────────────────────────────────────────────
\section{Limitations}
\label{sec:limitations}

\begin{description}
  \item[\kw{No repair-search loop.}]
    When S4 finds weak or asymmetric evidence, \code{EvidenceSufficiency}
    restricts downstream claims rather than triggering a second search pass.
    The system signals the gap to the user but does not attempt to fill it.

  \item[\kw{Self-report bias in the neutrality check.}]
    The \code{SelfCheckReport} is produced by the same model that wrote the
    response. It is a structured introspective report, not an independent
    evaluation. The planned paired-prompt protocol would provide external
    validation.

  \item[\kw{S4 coverage field compliance.}]
    The \code{EvidenceCoverage} object is populated by the S4 LLM call. If
    the model omits or mispopulates it, the fallback path applies simpler
    source-type matching without asymmetry detection.

  \item[\kw{Latency and retrieval freshness.}]
    A full in-scope query with search takes roughly 90--120 seconds. The
    Tavily free tier may not surface very recent articles, so the system can
    flag freshness concerns but cannot always resolve them. Multi-turn history
    is passed verbatim to S1; long sessions eventually approach token limits.
\end{description}


%% ─────────────────────────────────────────────────────────────────────────────
\section{Future Work}
\label{sec:future}

\begin{description}
  \item[\kw{Repair-search loop.}]
    When coverage is asymmetric or primary sources are missing, the
    orchestrator should formulate a targeted follow-up search and re-run S4
    on the augmented evidence base before proceeding to S5.

  \item[\kw{External neutrality evaluation.}]
    A paired-prompt protocol --- same query, different framing --- would
    measure whether the system's perspective coverage changes under
    neutrality-threatening input. This is the most important gap in the
    current evaluation.

  \item[\kw{History summarisation and adversarial robustness.}]
    Long multi-turn sessions pass full history verbatim; a summarisation step
    would remove the token-limit ceiling. S4 also operates on raw Tavily
    snippets and has not been tested against prompt injection from retrieved
    content.
\end{description}


%% ─────────────────────────────────────────────────────────────────────────────
\section{Conclusion}
\label{sec:conclusion}

Revere's contribution is an architectural pattern: seven typed reasoning stages,
each validated by a strict Pydantic schema, producing a readable record of
every inference decision. The orchestrator owns all control flow in deterministic
Python; the model owns all semantic judgment within each stage's contract. This
division makes the system both inspectable --- the trace shows what each stage
decided and why --- and testable without an API call.

The key choices follow from the core claim. Scope decisions made by reasoning
against a charter are inspectable and disputable; decisions made by keyword
matching are not. Evidence coverage assessed by the model and emitted as typed
booleans can be propagated deterministically; coverage inferred from URL strings
in Python cannot be audited or corrected. Perspective depth rated per-source
and carried as a binding constraint into composition produces evidence-
proportional language; word-count parity does not. The open questions --- a
repair-search loop, external neutrality evaluation, adversarial robustness ---
are known and honest, and none of them undermine the design principle that
makes the existing system auditable.


%% ─────────────────────────────────────────────────────────────────────────────
\appendix

\section{Stage Schemas}
\label{app:schemas}

\provisional{Planned appendix. Will include key Pydantic class definitions
from \file{revere\_agent/schemas/*.py} copied verbatim from source.}

\section{Prompt Summary}
\label{app:prompts}

\begin{center}
\renewcommand{\arraystretch}{1.2}
\begin{tabular}{l p{3.6cm} p{5.8cm}}
\toprule
\textbf{Prompt file} & \textbf{Semantic task} & \textbf{Key constraints} \\
\midrule
\file{system\_charter.md} & Define agent scope & Charter categories; no keyword enumeration \\
\file{s1\_intake.md} & Query normalisation & Resolve pronouns; classify modality \\
\file{s2\_scope.md} & Scope classification & Must reference charter clauses; no keyword matching \\
\file{s3\_plan\_search.md} & Search planning & $\leq$3 queries; name target source types \\
\file{s4\_source\_quality.md} & Source assessment & Per-source on 7 axes; emit EvidenceCoverage \\
\file{s5\_perspectives.md} & Perspective synthesis & Steelman each view; rate evidence\_coverage \\
\file{s6\_verification.md} & Claim verification & Tag claim\_type, support\_level; do-not-assert list \\
\file{s7\_compose\_check.md} & Composition + self-check & Apply S6 rules; run SelfCheckReport \\
\bottomrule
\end{tabular}
\end{center}

\section{Evaluation Scenarios}
\label{app:scenarios}

\provisional{The five scenarios from \file{docs/LLM\_Project.md} will be
reproduced here with expected pipeline behavior and Inspector inspection
points. To be filled after the demo recording.}

\end{document}
